\documentclass[12pt, a4paper, oneside]{report}

\usepackage[utf8]{inputenc}
\usepackage[czech]{babel}
\usepackage[T1]{fontenc}

\usepackage{lmodern}

\usepackage{graphicx}
\usepackage[hidelinks]{hyperref}
\usepackage{geometry}

\geometry{
    a4paper,
    left=2.5cm,
    right=2.5cm,
    top=3cm,
    bottom=3cm
}
\renewcommand{\baselinestretch}{1.4}

\begin{document}

\begin{titlepage}
    \begin{center}
        {DELTA – Střední škola informatiky a ekonomie, s.r.o.}\\
        Ke Kamenci 151, Pardubice \\

        \vfill

        {\Huge \textbf{Lifemeter}}\\
        {\Large Mobilní aplikace}\\

        \vfill

        \begin{flushleft}
            \textbf{Příjmení, jméno:} Zadrobílek Jan\\
            \textbf{Třída:} 4.B\\
            \textbf{SStudijní obor:} Informační technologie\\
            \textbf{Školní rok:} 2025/2026
        \end{flushleft}
    \end{center}
\end{titlepage}

\newpage
\thispagestyle{empty}
\mbox{}

\newpage
\thispagestyle{empty}
\noindent
Prohlašuji, že jsem maturitní projekt vypracoval(a) samostatně, výhradně s použitím uvedené literatury.\\

\noindent
V Pardubicích \number\day.~\number\month.~\number\year \hfill .................. \\

\newpage
\thispagestyle{empty}
\noindent
Rád bych upřímně poděkoval své vedoucí práce, Ing. Monice Borkovcové, Ph.D., za odborné vedení, cenné rady, trpělivost a čas, který mi věnovala při konzultacích a zpracování tohoto maturitního projektu.

\newpage
\thispagestyle{empty}
\section*{Resumé}
 [Stručný obsah práce v češtině - vyjádření hlavní myšlenky, problému a zvoleného řešení.]

\noindent \textbf{Klíčová slova:} [Klíčové slovo 1, Klíčové slovo 2, ... max 8 pojmů]


\section*{Summary (Resumé v angličtině)}
 [Brief summary of the thesis in English.]

\noindent \textbf{Keywords:} [Keyword 1, Keyword 2, ...]

\newpage
\thispagestyle{empty}
\tableofcontents

\newpage
\pagestyle{plain}

\chapter*{Úvod}
\addcontentsline{toc}{chapter}{Úvod}
[Vymezení problému, k čemu má práce sloužit, přehled doposud zkoumané problematiky.]

\chapter{Teoretická část}
 [Zmapování současného stavu, volba optimálních nástrojů, rešerše. Pozor na délku popisu technologií.]

\chapter{Metodika a vlastní řešení}
 [Postup práce, techniky. U SW projektů zde patří popis architektury, případně dokumentace API a UML diagramy.]

\chapter{Výsledky}
 [Stručně, jasně srozumitelně popsaná fakta bez úvah a komentářů.]

\chapter{Diskuse}
 [Zdůvodnění vybraného řešení, porovnání s jinými variantami. Zde nešetři místem.]

\chapter*{Závěr}
\addcontentsline{toc}{chapter}{Závěr}
[Celkové hodnocení, podmínky pro úspěšnou realizaci, stanovení dalšího možného rozšíření projektu.]

\newpage
\addcontentsline{toc}{chapter}{Literatura}
\begin{thebibliography}{99}
    \bibitem{vzor_kniha} PŘÍJMENÍ, J. \textit{Název knihy}. 1. vyd. Město: Nakladatelství, rok. ISBN.
    \bibitem{vzor_web} NÁZEV WEBU. \textit{Název článku} [online]. [Cit. 10.3.2026]. Dostupné z URL: http://...
\end{thebibliography}

\newpage
\addcontentsline{toc}{chapter}{Seznam obrázků, grafů a tabulek}
\listoffigures
\listoftables

% 10. Seznam příloh
\newpage
\chapter*{Seznam příloh}
\addcontentsline{toc}{chapter}{Seznam příloh}
\noindent Příloha 1: [Název první přílohy]\\
\noindent Příloha 2: [Název druhé přílohy]

% 11. Přílohy
\appendix
\chapter{Název první přílohy}
 [Zde vlož velké tabulky, diagramy, grafy, zdrojové kódy, které se nevešly do textu.]

\end{document}